\chapter{PIC behaviour and design decisions}
\label{ch:picbehav}

This chapter records the decisions the feature specification delegated, together
with what the RTL actually does. Each section names the requirement area from the
brief.

\section{Priority grouping}

Four bands exist. A band is a class of service rather than a fixed rank: its
urgency is a 2-bit value in \reg{BAND\_CONFIG}, so the four bands can be
reordered system-wide with a single register write while every source keeps its
own configuration.

Within a band, the 4-bit \sig{INTRA} field decides, higher being more urgent. If
those are equal too, the lowest source index wins. The two levels are separate
deliberately, because they change for different reasons: the band mapping is what
gets re-shuffled when the system changes operating mode, while the intra-band
value is the fine ordering inside a class.

Reconfiguring a source's band while an interrupt from it is already active cannot
corrupt nesting. The priority key is snapshotted onto the stack at the claim, so
the threshold in use is the value captured then. New pending evaluations of that
source use the new band immediately.

\section{Preemption and nesting}

Maximum depth is \reg{NEST\_MAX}, default 8, clamped by hardware to the range
[1,\,16]; 16 is the physical stack depth. Enforcement is by masking: once
\sig{depth} equals \reg{NEST\_MAX} the resolver offers nothing, so the CPU is
never shown an interrupt it could not legally claim. The limit therefore cannot
be exceeded by construction.

What is saved at a preemption is the source id and its priority key, pushed as a
pair onto the stack inside the controller. The key is what keeps the preemption
threshold stable; the id is what the end-of-interrupt pops. Architectural CPU
state --- \reg{mepc}, \reg{mcause}, \reg{mstatus} --- is the software's
responsibility under the standard RISC-V convention, so the controller stores
only what the controller itself needs.

The CPU signals completion with the dedicated \sig{cpu\_irq\_eoi} pulse.

A dedicated pulse is what makes nesting possible. The CPU's \sig{cpu\_in\_trap}
level says only that a trap is in progress; it cannot express \emph{which}
nesting level has finished, so one pulse per interrupt-returning \sig{MRET} is
required to pop exactly one. \sig{cpu\_in\_trap} is still driven, for
observability, but the PIC does not use it.

A preemption blocked by the depth limit is visible rather than silent:
\reg{INT\_STATUS.OVF} and \reg{NEST\_STATUS.OVF} are set. Nothing is lost --- the
blocked source stays pending and is offered as soon as the depth drops.

\section{Spurious interrupt detection}

A source may deassert between being offered and the CPU's claim. The condition is
detectable at exactly one point: the claim itself. When \sig{cpu\_irq\_ack}
arrives, the controller re-checks the request of the source it had offered; if
that request is gone, the claim is spurious.

\begin{center}
\ttfamily\small
claim\_spur = claim\_push \&\& !req[claimed\_id];
\end{center}

The response is status flags, not a discard and not a dedicated vector:
\reg{SRCx\_STATUS.SPUR}, the sticky \reg{SPURIOUS\_LOG} bit, and
\reg{INT\_STATUS.SPUR}. The claim is \emph{still} pushed onto the nesting stack.

That last point looks wrong at first reading and is the deliberate choice. A
vectored CPU has already committed to entering the handler by the time the PIC
learns the request vanished. If the PIC skipped the push, that handler's
end-of-interrupt would pop somebody else's level and corrupt the stack. Keeping
\sig{ack} and \sig{eoi} balanced matters more than keeping the depth counter
tidy. The handler reads the flag and exits early. A dedicated spurious vector was
rejected because it would require a separate CPU entry path for a case that costs
a two-instruction check.

Distinguishing a spurious interrupt from a legitimately fast-clearing source is a
question of \emph{when} the request disappears, not how quickly. A source still
asserted at the claim and cleared later inside its own handler is the normal case
--- that is how every peripheral in this system behaves --- and is never flagged.
Only a request already gone at the moment of the ack is spurious. Edge-triggered
sources are immune by construction, since the latch holds the request until it is
consumed.

Logging is at two granularities: \reg{SPURIOUS\_LOG} identifies \emph{which}
source misbehaved and survives until software clears it, while
\reg{INT\_STATUS.SPUR} is the cheap global poll.

\section{Deadline-aware escalation}
\label{sec:deadline}

Each source has a 16-bit deadline in \reg{SRCx\_CONFIG.DEADLINE}, in the same
word as its band and priority, so one write configures a source completely. A
value of zero disables the deadline, which is also the reset state.

Elapsed time is measured by a per-source counter that runs while the slot is a
pending request awaiting service, and is cleared on claim or on return to idle.
It is not free-running: it measures exactly the interval the deadline is about.

\begin{table}[H]
\centering
\caption{Deadline counter update (\file{hdl/pic.v})}
\label{tab:ddlcnt}
\begin{tabular}{@{}lL{0.60\linewidth}@{}}
\toprule
\textbf{Condition} & \textbf{Action on \sig{ddl\_cnt[s]}} \\
\midrule
not counting        & cleared --- the slot is idle or has no deadline \\
escalation fires    & cleared --- re-arms the window \\
deadline not yet hit& increments \\
\bottomrule
\end{tabular}
\end{table}

Escalation means one of two things, selected by \reg{ESCALATION\_CFG.MODE}. Jump
mode sets the effective band to the configurable \sig{TARGET} --- one decisive
promotion regardless of where the source started. Bump mode moves it one band
more urgent, stopping at band 0 --- a gentler, proportional escalation. Only the
\emph{effective} band changes; the configured band in \reg{SRCx\_CONFIG} is left
untouched, so the escalation is transparent and reversible.

Repetition is optional. With \sig{MULTI} = 0 a source escalates at most once per
waiting episode. With \sig{MULTI} = 1 the counter re-arms after each escalation,
so a source that keeps waiting keeps climbing; combined with bump mode this gives
a source that ages upward one band per deadline period.

The escalated priority is dropped at the claim: from the cycle the source goes
into service it is no longer a pending request awaiting service, so the
effective band reverts to the configured band and the deadline counter resets.
Escalation therefore never accumulates across events --- each request starts
from the priority software actually programmed.
\reg{SRCx\_STATUS.ESC} clears with it, while the global \reg{INT\_STATUS.ESC}
stays sticky so the event itself is not lost.

The promotion takes effect on the cycle the deadline is missed, not on the one
after it: the escalated band is already in the priority key the resolver
evaluates on that same cycle.

\section{Software-triggered interrupts}

A software interrupt is raised by a single write to \reg{SRCx\_SW\_TRIG} with bit
0 set and the key \texttt{0xA5A5} in bits [31:16]. One write, not a sequence, so
there is no partially armed state to get stuck in. The key is what makes it safe:
a stray or corrupted single-bit write does nothing.

Clearing works both ways. Writing bit 0 = 0 clears it explicitly, with no key
required, and a claim also self-clears it exactly like an edge-triggered hardware
source, so a handler need not remember to tidy up after itself.

If a software trigger lands on a channel that already has a pending hardware
request, the two merge and the slot simply stays pending --- the request is
\sig{(hardware | software) \& enable}, so there is one request per slot rather
than two competing ones. The consequence is worth stating: a claim consumes the
software bit, but if the hardware line is still asserted the slot remains pending
and will be offered again. That is correct level-triggered behaviour, not a lost
or duplicated interrupt.

A software request is subject to deadline escalation exactly like any other, with
no special case in the RTL. That is the point of merging at the slot: once inside
the priority pipeline a software interrupt is genuinely indistinguishable from a
hardware one.

\section{Reset behaviour}

Every configuration and status register resets as listed in
Table~\ref{tab:picmap}, with the two documented exceptions. The nesting depth
resets to zero, the active, escalated, spurious and edge-pending masks all reset
to zero, and the output pair \sig{cpu\_irq} / \sig{cpu\_irq\_vec} resets to
\{0,\,0\}.

The reset is asynchronous and active low, matching the CPU. Reset assertion is
therefore visible at the instant it occurs rather than at the next clock edge;
release takes effect on the following rising edge.
